\documentclass[11pt]{article}
\usepackage[T1]{fontenc}
\usepackage{lmodern,microtype,amsmath,amssymb,amsthm,booktabs,array}
\usepackage[margin=0.94in]{geometry}
\usepackage[hidelinks]{hyperref}
\newtheorem{proposition}{Proposition}
\newtheorem{lemma}{Lemma}
\newtheorem{theorem}{Theorem}
\newcommand{\ab}{a_{\rm b}}
\newcommand{\ac}{a_{\rm c}}
\hypersetup{pdftitle={A linear guard and a constant-width Gaussian solve},
pdfauthor={William Porter (AI-assisted)}}
\title{A linear guard and a constant-width Gaussian solve\newline
\large Downstream assembly refinements for integer multiplication}
\author{William Porter\thanks{Prepared with AI assistance (Claude agents).
This note builds on odepoint's \emph{intmultspam}, a fork of Douglas Colkitt's
\emph{integer-mult-bounds}, and on OpenAI's paper \emph{Integer multiplication
below $n\log n$}. It has no independent review or endorsement from their
authors.}}
\date{October 8, 2026}
\begin{document}
\maketitle
\begin{abstract}
These notes change two downstream components of the conditional
integer-multiplication bound, leaving both finite networks untouched.
First, the coefficient guard of a simultaneous butterfly layer is bounded
linearly in the number of axes, by charging each completed recursive child
with its exact kernel instead of its internal dependency depth. The guard
condition $\epsilon C_1<1$ becomes $\epsilon<1$. Second, the Gaussian
resampling system is solved by a precomputed banded LU factorization with a
complete fixed-point error analysis. A diagonal similarity makes the system
strongly diagonally dominant for every separation $\theta$, so the Gaussian
width can be the constant $\alpha=2$. The Gaussian row of the cost table
drops from power $\frac34+\delta+\frac54\epsilon$ to
$\frac12+\delta+\epsilon$. For certified savings $\ab,\ac$ the parameter
supremum moves from $Q/(5+4Q)$ to $Q/(2+2Q)$, $Q=\min(\ab,\ac)$. With the
current certified savings, exact rational parameters give
\[
 T(n)=O\!\left(n(\log n)^{1-\kappa}\right),\qquad
 \kappa=\frac{296461013}{2\cdot10^{17}}=1.4823050650\cdot10^{-9}>2^{-30},
\]
about $2.5$ times the previous conditional witness $5929220328/10^{19}$.
The result is conditional on the same upstream interfaces as before, plus
the new written proofs. The exact checks support those proofs; they do not
independently establish the complete upstream theorem.
\end{abstract}

\section{Baseline, dependencies, and the change}
The baseline is odepoint's \emph{intmultspam}~\cite{fork}, revision
\texttt{8e1701f}, a fork of Douglas Colkitt's
\emph{integer-mult-bounds}~\cite{baseline}. Its conditional witness is
$\kappa=5929220328/10^{19}>2^{-31}$, which is the exact parameter supremum
$\ab/(5+4\ab)$ of its assembly for the certified bit saving
$\ab=29646101715649/10^{22}$. Both rest on OpenAI's \emph{Integer
multiplication below $n\log n$}~\cite{upstream}, pinned at
\begin{center}\texttt{adc7f1241b42e322a6451854ab7e4b4c146bf78a}.\end{center}
Numbered citations to Harvey and van der Hoeven~\cite{HarveyHoeven2021}
refer to its author-hosted manuscript, as in the pinned source.

We retain every interface that the baseline assumes: the pinned stream
operations, bit transfer and complex phase transfer, the paired bit network
and the pair-star complex network with their certified savings
$\ab=29646101715649/10^{22}$ and $\ac=4/10^9$, nonadjacent routing,
compact dirty controls with their reservations and local repair, synthetic
transforms, prime selection, the Gaussian identity and its forward and
diagonal maps, and the final rounding. Two components change.
\begin{enumerate}
\item The width guard of the simultaneous layer (Section~\ref{sec:note-guard}).
The retained depth recurrence $A(e)\le sA(e/m)+E$ gives widths
$p+O(d^{C_1})$ with $C_1=5-4\beta+\zeta$. We prove widths $p+O(d)$.
\item The solve of the square Gaussian system
(Section~\ref{sec:note-resampling}). The retained Neumann series costs
$O(tp^{5/2+\delta}/(\alpha^3\theta))$ on a line of length $t$ and is used
under the hypothesis $\theta>p/\alpha^4$, so that it does not dominate the
$O(tp^{3/2+\delta}\alpha)$ forward map~\cite[Section~4.4.2]{HarveyHoeven2021}.
The retained separations are $\theta_i\asymp1/d$, so $\alpha\asymp(pd)^{1/4}$,
and the Gaussian row has power $\frac34+\delta+\frac54\epsilon$: one factor
$d$ for the axes and $\alpha\asymp p^{(1+\epsilon)/4}$. Harvey and van der
Hoeven remark that a precomputed LU decomposition ``allows us to relax the
condition $\theta>p/\alpha^4$ to merely $\theta>1/\alpha^2$'', and that they
did not use it ``because the error analysis is considerably more intricate
than for the iterative method''~\cite[Section~4.4.2, pp.~34--35]{HarveyHoeven2021}.
With $\theta_i\asymp1/d$ that condition would still force
$\alpha\asymp\sqrt d$. We prove a complete version that needs neither
condition, for every integer $\alpha\ge2$, and use $\alpha=2$.
\end{enumerate}

\begin{center}
\begin{tabular}{@{}lll@{}}
\toprule
& Baseline & This note\\
\midrule
Coefficient guard & $\epsilon(5-4\beta+\zeta)<1$ & $\epsilon<1$\\
Gaussian width $\alpha$ & $\lceil(32db)^{1/4}\rceil$ & $2$\\
Gaussian row power & $\frac34+\delta+\frac54\epsilon$ & $\frac12+\delta+\epsilon$\\
Source scale $\gamma$ & $O(p^{1/2+3\epsilon/2})$ & $O(d^2)=O(p^{2\epsilon})$\\
Parameter supremum & $Q/(5+4Q)$ & $Q/(2+2Q)$\\
Witness $\kappa$ & $5.929220328\cdot10^{-10}$ & $1.4823050650\cdot10^{-9}$\\
\bottomrule
\end{tabular}
\end{center}

\section{Coefficient widths of the simultaneous layer}\label{sec:note-guard}
The simultaneous layer applies $C^{\otimes D}$, with
$C=aI+bX$, $a=(1+i)/2$, $b=(1-i)/2$, to $D\le d$ selected address bits, and
converts it to the normalized butterfly layer by
$H_0^{\otimes D}=b^D\mathcal S_DC^{\otimes D}\mathcal S_D$, where
$\mathcal S_D$ multiplies each record by $i$ to the number of selected bits
equal to one. Its recursion splits a group of $e=mf$ axes into $m$ slots,
applies the finite pair-star network with $W$ roles and $s$ child calls,
stops below $d^\beta$ axes, and uses individual kernels at the leaves and on
at most $d$ reserved axes. The node count (PS11a)--(PS11b) cited below is
the one in \texttt{notes/complex-pair-star-construction.tex}, which the
pair-star manuscript patch \texttt{patches/complex-pair-star-31.patch}
inserts into Section~3: it gives
$D_{\rm scalar}=3v^2(12R_{\rm side}+4vh+10v)+8W$ for the scalar circuit and
its endpoint and bank corrections, and $D_{\rm node}=D_{\rm scalar}+4s$,
where $4s$ pays for the inverse-child wrappers. The following subsection
replaces the guard subsection of the retained layer proof.

\subsection{A linear coefficient guard}
\label{sec:linear-guard}
The accelerated layer is evaluated exactly until the entire normalized
butterfly layer is complete, and its intermediate operators need not be
contractions. We bound the denominators and moduli of every intermediate
coefficient by a quantity that is \emph{linear} in $d$. The bound replaces the
dependency-depth recurrence $A(e)\le sA(e/m)+E$, which charges each of the
$s$ consecutive children of a node with the whole internal depth of that
child, and therefore grows like a fixed power of $d$. Here a \emph{completed}
child is charged only for its exact action, which is a tensor power of the
fixed kernel $C$.

Use the complex network's own constants $m\ge3$, $W$ and $s$, independently
of the bit network used by swaps. One network invocation is the computation
at one internal node: its scalar gates, endpoint and bank corrections,
inverse-child wrappers, binary changes of basis, row splitting and
reassembly, and exactly $s$ child calls. Let $D_{\rm node}$ bound the number
of elementary coefficient operations of one invocation outside its child
calls; an elementary operation is a complex addition, subtraction, halving,
sign change or multiplication by $i$, applied coefficientwise to aligned
records. For the pair-star network the finite count (PS11a)--(PS11b) of the
pair-star construction gives
\[
 D_{\rm node}<E=64(W+m+1)^3 .
\]
Binary changes of basis, compact-control swaps and rotations, exceptional
repair, row padding, parking on the stack, copying and bank exchanges
permute or copy complete coefficient encodings, or write zero rows; they
perform no coefficient arithmetic.

For an integer $P\ge0$ and real $M>0$, call a finite family of complex numbers
\emph{$(P,M)$-bounded} if all its members lie in $2^{-P}\mathbb Z[i]$ and have
modulus at most $M$. At a given moment, the \emph{state family} of a call is
the family of all coefficient values in that call's streams, including
scratch roles, padded zero rows and parked streams.

\begin{lemma}\label{lem:guard-steps}
If the state family is $(P,M)$-bounded and one elementary operation is applied
coefficientwise, the new state family is $(P+1,2M)$-bounded. Permutations and
copies of complete coefficient encodings, and insertion of zero rows, preserve
$(P,M)$-boundedness.
\end{lemma}
\begin{proof}
The set $2^{-P}\mathbb Z[i]$ is closed under sums, differences, sign changes
and multiplication by $i$, and $|a\pm b|\le2M$; halving maps it into
$2^{-P-1}\mathbb Z[i]$ and does not increase moduli. Values that are not
changed keep their bounds, and a $(P,M)$-bounded family is
$(P+1,2M)$-bounded. Movement creates no new values, and $0$ lies in every
such family.
\end{proof}

\begin{lemma}\label{lem:guard-kernels}
Let $f\ge1$ and let $\mathcal K$ be $C^{\otimes f}$ or
$(C^{-1})^{\otimes f}$, acting on $f$ selected address bits of a stream,
coefficientwise and as the identity on all other address bits. If the values
of the stream are $(P,M)$-bounded, then those of $\mathcal K$ applied to the
stream are $(P+\lceil f/2\rceil,2^{f/2}M)$-bounded.
\end{lemma}
\begin{proof}
Put $a=(1+i)/2$ and $b=(1-i)/2=-ia$. For $x,y\in\mathbb F_2^f$ with $k$ the
Hamming weight of $x+y$, the entry of $C^{\otimes f}$ at $(x,y)$ is
$a^{f-k}b^k=(-i)^ka^f=(-i)^k(1+i)^f2^{-f}$, and the entry of
$(C^{-1})^{\otimes f}$, with $C^{-1}=bI+aX$, is
$b^{f-k}a^k=(-i)^{f-k}a^f$. Since $(1+i)^2=2i$, each entry is a unit times
$(1+i)^{f\bmod2}2^{-\lceil f/2\rceil}$: it lies in
$2^{-\lceil f/2\rceil}\mathbb Z[i]$ and has modulus $2^{-f/2}$. An output
coefficient is a $\mathbb Z[i]$-combination of $2^f$ input coefficients with
such entries. Hence it lies in $2^{-P-\lceil f/2\rceil}\mathbb Z[i]$, and the
maximum absolute row sum $2^f\cdot2^{-f/2}$ bounds its modulus by
$2^{f/2}M$.
\end{proof}

For a node with $e$ active axes, an integer $\Gamma$ is a \emph{valid growth}
if, whenever the node's input family is $(P,M)$-bounded, every value occurring
during its execution, including inside all descendants, is
$(P+\Gamma,2^\Gamma M)$-bounded.

\begin{lemma}\label{lem:guard-node}
A leaf with $e<d^\beta$ active axes has valid growth $8e$. If every child of an
internal node with $e=mf$ active axes has valid growth $\Gamma_f$, then
\[
 \Gamma_e=D_{\rm node}+s\lceil f/2\rceil+\Gamma_f
\]
is a valid growth for the node. A completed node with $e$ active axes maps a
$(P,M)$-bounded input family to a $(P+\lceil e/2\rceil,2^{e/2}M)$-bounded
output family.
\end{lemma}
\begin{proof}
A leaf applies $e$ individual kernels, each with at most eight elementary
operations; apply Lemma~\ref{lem:guard-steps} at most $8e$ times.

At an internal node, apart from movement, the invocation performs at most
$D_{\rm node}$ elementary operations and exactly $s$ child calls on $f$ axes.
A completed forward child acts on its role stream as $C^{\otimes f}$ on the
selected bits of its slot and as the identity elsewhere. An inverse child is
one forward child between the unit wrappers $Z^{\otimes f}$ and the scalar
$(-i)^f$, so its completed action is $(C^{-1})^{\otimes f}$; the wrappers are
among the counted elementary operations. By Lemmas~\ref{lem:guard-steps}
and~\ref{lem:guard-kernels}, just before the $j$-th child call the state
family is $(P+a_j,2^{a_j}M)$-bounded with
$a_j\le D_{\rm node}+(j-1)\lceil f/2\rceil$, because $f/2\le\lceil f/2\rceil$.
The child's input is a subfamily and the parent's parked values are unchanged
during the call, so every value during the call is within
$P+a_j+\Gamma_f$ and $2^{a_j+\Gamma_f}M$. After all steps the bound is
$D_{\rm node}+s\lceil f/2\rceil$, which proves the second claim.

Every completed invocation applies $C^{\otimes e}$ to every row of every
role, including roles whose initial scratch values were arbitrary, and
padded zero rows remain zero. Lemma~\ref{lem:guard-kernels} with $f=e$ gives
the last claim.
\end{proof}

\begin{proposition}[Linear guard]\label{prop:linear-guard}
Put
\[
 C_0^*=E+2s+m+26,\qquad \Delta=C_0^*d .
\]
For every input of the simultaneous layer whose coefficients are disk values
in $2^{-p}\mathbb Z[i]$, and for every $1\le D\le d$ and fixed rational
$0<\beta<1$, every intermediate coefficient lies in
$2^{-(p+\Delta)}\mathbb Z[i]$ and has modulus at most $2^\Delta$. Thus $p+\Delta$
fractional bits and $\Delta+3$ integer and sign bits per real component
represent every intermediate value exactly. If $d\le b_dp^\epsilon$ with a
fixed $0<\epsilon<1$, then $\Delta\le p$ beyond a fixed cutoff, so these
widths are $O(p)$.
\end{proposition}
\begin{proof}
The input family is $(p,1)$-bounded. The layer first applies at most $d$
individual kernels on its reserved row, front and back axes, or on all $D\le d$
axes in the small-$D$ fallback; by Lemma~\ref{lem:guard-steps} all values stay
within $p+8d$ and $2^{8d}$.

The remaining active axes form $n\le(m-1)(1+\lfloor\log_md\rfloor)\le(m-1)d$
base-$m$ pieces of sizes $e_1,\ldots,e_n$ with $\sum_je_j\le d$, executed one
after another. For a piece of size $e_0$ with internal sizes
$e_i=e_0/m^i$ and leaf size $e_k<d^\beta$, Lemma~\ref{lem:guard-node} gives,
using $\lceil x\rceil\le x+1$, $k\le\lceil\log_md\rceil\le d$ and $e_0\le d$,
\[
 \Gamma_{e_0}\le\sum_{i=1}^k\bigl(D_{\rm node}+s\lceil e_i/2\rceil\bigr)+8e_k
 \le d(D_{\rm node}+s)+\frac{sd}{2(m-1)}+8d .
\]
When piece $j$ completes, the family bound has grown by at most
$\lceil e_j/2\rceil$ in the denominator exponent and $e_j/2$ in the binary
logarithm of the modulus, by the last claim of Lemma~\ref{lem:guard-node}. Hence during piece $j$ every value is within
$p+8d+\sum_{i<j}\lceil e_i/2\rceil+\Gamma_{e_j}$, and after all pieces the
family is within $p+8d+d/2+n$.

The two outer diagonals $\mathcal S_D$ are unit multiplications. The scalar
$b^D=(-i/2)^tb^{D\bmod2}$, $t=\lfloor D/2\rfloor$, consumes $t\le d/2$
fractional bits through the shift by $t$, which only lowers moduli, and at most
nine further elementary operations. Altogether every intermediate value is
within
\[
 d\Bigl(D_{\rm node}+s+\frac{s}{2(m-1)}+m+16\Bigr)+9
 \le d(D_{\rm node}+2s+m+25)\le C_0^*d,
\]
using $s/(2(m-1))\le s$, $9\le9d$ and $D_{\rm node}<E$. The modulus exponent
obeys the same bound, because each step above increases it by at most the
amount charged to the denominator exponent. A value in
$2^{-(p+\Delta)}\mathbb Z[i]$ of modulus at most $2^\Delta$ has real and
imaginary numerators of absolute value at most $2^{p+2\Delta}$, so $p+\Delta$
fractional and $\Delta+3$ integer and sign bits hold it exactly. Finally
$C_0^*d\le C_0^*b_dp^\epsilon\le p$ for all large $p$ when $\epsilon<1$.
\end{proof}

All compact-control operations, exceptional repair, padding and parking were
counted as movement. The individually processed axes are disjoint and number
at most $d$, which is the same $8d$ allowance used above. Thus the condition
$\epsilon C_1<1$ of the dependency-depth guard is replaced by $\epsilon<1$,
and the stopping exponent $\beta$ no longer enters any width condition.


\section{Gaussian resampling with a constant width}\label{sec:note-resampling}
Here $\mathbb D_p$ is the disk grid $2^{-p}\mathbb Z[i]\cap\{|z|\le1\}$, vectors
carry the maximum coordinate norm, and an approximation $\widetilde L$ of a
linear contraction $L$ has error $E$ if $2^p\|\widetilde Lu-Lu\|<E$ for every
input on the grid. The maps $\mathsf S,\mathsf T,C,\mathsf D$ are those
of~\cite[Section~4]{HarveyHoeven2021}: for $0\le k<t$,
\[
 (\mathsf Su)_k=\frac1\alpha\sum_{j\in\mathbb Z}e^{-\pi\alpha^{-2}(j-sk/t)^2}u_j,
 \qquad
 (\mathsf Tu)_k=\sum_{j\in\mathbb Z}e^{-\pi\alpha^2(tj/s-k)^2}u_j,
\]
$C$ selects the coordinates $q_\ell$ and
$\mathsf D={\rm diag}(e^{\pi\alpha^2\beta_\ell^2})$, so that $N=C\mathsf T\mathsf D$.

\subsection{A precomputed banded solve}
\label{sec:banded-solve}
The cited one-dimensional resampling solves its square Gaussian system by a
Neumann series~\cite[Lemma~4.12]{HarveyHoeven2021}, which needs
$\theta>p/\alpha^4$. Its Section~4.4.2 remarks that a precomputed LU
decomposition would only need $\theta>1/\alpha^2$, but the authors did not use
it ``because the error analysis is considerably more intricate''. We prove a
precomputed banded solve with a complete fixed-point error analysis. With a
diagonal rescaling it works for \emph{every} $\theta$ and every
$\alpha\ge2$; it is therefore used below with the constant width $\alpha=2$.

Throughout this subsection let $s<t<2s$ be coprime positive integers, put
$\theta=t/s-1\in(0,1)$, and let $\alpha\ge2$. For $\ell\in\mathbb Z$ write
\[
 x_\ell=\frac{t\ell}s,\qquad q_\ell=\Bigl\lfloor x_\ell+\frac12\Bigr\rfloor,
 \qquad \beta_\ell=x_\ell-q_\ell\in\Bigl[-\frac12,\frac12\Bigr),
 \qquad \Delta q_\ell=q_{\ell+1}-q_\ell,
\]
and, for $h\in\mathbb Z$,
\[
 Q_h(\ell)=\Bigl(\frac{th}s+\beta_\ell\Bigr)^2-\beta_{\ell+h}^2 .
\]
For $0\le\ell<s$ these are the selected coordinates $q_\ell$ and fractional
parts $\beta_\ell$ of the cited Gaussian system, and $\beta_{\ell+s}=\beta_\ell$.
The square matrix $N=C\mathsf T\mathsf D$ of the cited
Section~4.2, written with indices modulo $s$, has entries
\[
 N_{\ell,j}=\sum_{h\equiv j-\ell\ ({\rm mod}\ s)}e^{-\pi\alpha^2Q_h(\ell)} ,
\]
and $N=I+E$, where $E$ collects all terms with $h\ne0$; every such term is a
positive real number. For matrices, $\|\cdot\|$ is the maximum absolute row
sum, induced by the maximum coordinate norm.

\begin{lemma}[Diagonally dominant LU factors]\label{lem:dominant-lu}
Let $M$ be a real $n\times n$ matrix with
$|M_{ii}-1|+\sum_{j\ne i}|M_{ij}|\le\rho$ for every $i$, where
$0\le\rho\le\frac1{10}$. Gaussian elimination without pivoting runs to
completion and gives $M=LU$, with $L$ unit lower triangular and $U$ upper
triangular. Every Schur complement $S$ in the elimination satisfies
$S_{ii}\in[1-\rho,1+\rho]$ and $\sum_{j\ne i}|S_{ij}|\le\rho$ in each row.
Moreover, for every $i$,
\[
 U_{ii}\in[1-\rho,1+\rho],\qquad \sum_{j>i}|U_{ij}|\le\rho,\qquad
 \sum_{j<i}|L_{ij}|\le\frac{\rho}{1-2\rho}\le\frac18,
\]
and, with $D_U$ the diagonal of $U$,
\[
 \|L^{-1}\|\le\tfrac87,\qquad\|(D_U^{-1}U)^{-1}\|\le\tfrac98,\qquad
 \|U^{-1}\|\le\tfrac54,\qquad\|M^{-1}\|\le\tfrac{10}9 .
\]
\end{lemma}

\begin{proof}
For a matrix $X$ with positive diagonal write
$\omega_i(X)=\sum_{j\ne i}|X_{ij}|$, $\mu_i(X)=X_{ii}-\omega_i(X)$ and
$\tau_i(X)=X_{ii}+\omega_i(X)$. Initially $\mu_i(M)\ge1-\rho>0$ and
$\tau_i(M)\le1+\rho$. One elimination step with pivot $1$ forms
$X'_{ij}=X_{ij}-X_{i1}X_{1j}/X_{11}$ for $i,j\ge2$. Then
\[
 \omega_i(X')\le\omega_i(X)-|X_{i1}|+\frac{|X_{i1}|}{X_{11}}
 \bigl(\omega_1(X)-|X_{1i}|\bigr),\qquad
 |X'_{ii}-X_{ii}|\le\frac{|X_{i1}||X_{1i}|}{X_{11}},
\]
hence $\mu_i(X')\ge\mu_i(X)+|X_{i1}|\mu_1(X)/X_{11}\ge\mu_i(X)$ and
$\tau_i(X')\le\tau_i(X)-|X_{i1}|\mu_1(X)/X_{11}\le\tau_i(X)$. By induction
every Schur complement $S$ has positive pivots,
$\mu_i(S)\ge1-\rho$ and $\tau_i(S)\le1+\rho$; thus
$S_{ii}\in[1-\rho,1+\rho]$ and $\omega_i(S)=(\tau_i(S)-\mu_i(S))/2\le\rho$.
Row $k$ of $U$ is row $k$ of the $k$-th Schur complement $S^{(k)}$, which gives
the bounds on $U$.

For $L$, fix $i$ and put $\omega^{(k)}=\sum_{j\ge k,\,j\ne i}|S^{(k)}_{ij}|$.
Eliminating pivot $k<i$ removes $|S^{(k)}_{ik}|$ from row $i$ and adds at most
$(|S^{(k)}_{ik}|/S^{(k)}_{kk})\,\omega_k(S^{(k)})\le|S^{(k)}_{ik}|\rho/(1-\rho)$
to its other off-diagonal entries. Hence
$|S^{(k)}_{ik}|(1-2\rho)/(1-\rho)\le\omega^{(k)}-\omega^{(k+1)}$, and summing
over $k<i$ gives $\sum_{k<i}|S^{(k)}_{ik}|\le\rho(1-\rho)/(1-2\rho)$. As
$L_{ik}=S^{(k)}_{ik}/S^{(k)}_{kk}$ and $S^{(k)}_{kk}\ge1-\rho$, the row bound on
$L$ follows. Finally $L-I$, $D_U^{-1}U-I$ and $M-I$ have norms at most
$\frac18$, $\frac{\rho}{1-\rho}\le\frac19$ and $\frac1{10}$, and
$\|D_U^{-1}\|\le\frac{10}9$; Neumann series give the four norm bounds.
\end{proof}

The matrix $N$ itself need not be diagonally dominant: when
$\beta_{\ell+1}$ is close to $-\frac12$, the entry $N_{\ell,\ell+1}$ is close to
one. A diagonal similarity removes this. Let $\psi$ be the $1$-periodic
function with
\[
 \theta\,\psi(x)=\begin{cases}
 x\bigl(x+\frac12\bigr),&-\frac12\le x\le-\frac14,\\[2pt]
 -\frac1{16},&-\frac14\le x\le\frac14,\\[2pt]
 x\bigl(x-\frac12\bigr),&\frac14\le x\le\frac12,
 \end{cases}
\]
and put
\[
 \Phi=\frac1{16\theta},\qquad \varphi_\ell=\psi(\beta_\ell)+\Phi\in[0,\Phi],
 \qquad \Lambda={\rm diag}\bigl(e^{\pi\alpha^2\varphi_\ell}\bigr)_{\ell\in\mathbb Z/s},
 \qquad R=e^{\pi\alpha^2\Phi}.
\]
Thus every diagonal entry of $\Lambda$ lies in $[1,R]$.

\begin{lemma}[Scaled dominance]\label{lem:scaled-dominance}
For every $\ell$ and $h\ne0$ the scaled exponent
$\mathcal E_{\ell,h}=Q_h(\ell)-(\varphi_{\ell+h}-\varphi_\ell)$ satisfies
\[
 \mathcal E_{\ell,\pm1}\ge\tfrac12+\theta,\qquad
 \mathcal E_{\ell,h}\ge|h|\bigl(|h|-\tfrac32\bigr)\quad(|h|\ge2).
\]
Consequently $\Lambda^{-1}N\Lambda=I+\Lambda^{-1}E\Lambda$ with
\[
 \|\Lambda^{-1}E\Lambda\|\le\rho_\alpha:=2e^{-\pi\alpha^2/2}+2.01e^{-\pi\alpha^2}
 \le\rho_2<0.0038,
\]
$N$ is invertible with $\|N^{-1}\|\le R/(1-\rho_2)\le1.004R$, and for every
integer $m\ge1$ the terms with $|h|>m$ contribute less than
$3e^{-\pi\alpha^2m^2}$ to each row of $\Lambda^{-1}E\Lambda$.
\end{lemma}

\begin{proof}
(a) Since $x_{\ell+1}-x_\ell=1+\theta\in(1,2)$ and $|\beta|\le\frac12$, we have
$\Delta q_\ell\in\{1,2\}$ and $\beta_{\ell+1}=\beta_\ell+\theta+1-\Delta q_\ell$.

(b) As $t/s+\beta_\ell=x_{\ell+1}-q_\ell=\beta_{\ell+1}+\Delta q_\ell$ and
$-t/s+\beta_{\ell+1}=x_\ell-q_{\ell+1}=\beta_\ell-\Delta q_\ell$,
\[
 Q_1(\ell)=\Delta q_\ell(\Delta q_\ell+2\beta_{\ell+1}),\qquad
 Q_{-1}(\ell+1)=\Delta q_\ell(\Delta q_\ell-2\beta_\ell).
\]
(c) For $h\ne0$, $|th/s+\beta_\ell|\ge(t/s)|h|-\frac12>|h|-\frac12\ge\frac12$,
so $Q_h(\ell)>(|h|-\frac12)^2-\frac14=|h|(|h|-1)\ge0$, as in the proofs
of~\cite[Lemmas~4.6 and~4.11]{HarveyHoeven2021}.

(d) We have $\psi(x)=\theta^{-1}\int_{-1/2}^xg$, where $g(y)=2y+\frac12$ on
$[-\frac12,-\frac14]$, $g=0$ on $[-\frac14,\frac14]$ and $g(y)=2y-\frac12$ on
$[\frac14,\frac12]$. On $[-\frac12,\frac12]$ one has
$2y-\frac12\le g(y)\le2y+\frac12$ and $|g(y)|\le\frac12$, and
$\int_{-1/2}^{1/2}g=0$; so $\psi$ is continuous and $1$-periodic, and
$\theta\psi$ takes values in $[-\frac1{16},0]$.

(e) In both cases of (a), periodicity gives
$\varphi_{\ell+1}-\varphi_\ell=\psi(\beta_\ell+\theta)-\psi(\beta_\ell)
=\theta^{-1}\int_{\beta_\ell}^{\beta_\ell+\theta}\tilde g$, where $\tilde g$ is
the periodic extension of $g$. Hence
$|\varphi_{\ell+1}-\varphi_\ell|\le\frac12$ for every $\ell$. If
$\Delta q_\ell=1$, then $[\beta_\ell,\beta_\ell+\theta]\subset[-\frac12,\frac12)$,
and averaging the bounds of (d) gives
\[
 2\beta_\ell+\theta-\tfrac12\le\varphi_{\ell+1}-\varphi_\ell\le2\beta_\ell+\theta+\tfrac12 .
\]
(f) If $\Delta q_\ell=1$, then by (b)
$\mathcal E_{\ell,1}=1+2\beta_\ell+2\theta-(\varphi_{\ell+1}-\varphi_\ell)\ge\frac12+\theta$
and $\mathcal E_{\ell+1,-1}=1-2\beta_\ell+(\varphi_{\ell+1}-\varphi_\ell)\ge\frac12+\theta$.
If $\Delta q_\ell=2$, then $Q_1(\ell)\ge2$ and $Q_{-1}(\ell+1)\ge2$, and both
scaled exponents are at least $\frac32$.

(g) For $|h|\ge2$, (e) gives $|\varphi_{\ell+h}-\varphi_\ell|\le|h|/2$, and (c)
gives $\mathcal E_{\ell,h}\ge|h|(|h|-1)-|h|/2$.

(h) Row $\ell$ of $\Lambda^{-1}E\Lambda$ is the sum, over all $h\ne0$, of the
positive terms $e^{-\pi\alpha^2\mathcal E_{\ell,h}}$; the terms with
$h\equiv0\pmod s$ lie on the diagonal. By (f) and (g) the total is at most
$2e^{-\pi\alpha^2/2}+2\sum_{k\ge2}e^{-\pi\alpha^2k(k-3/2)}$, and successive
terms of the last sum decrease by at least $e^{-14\pi}$ when $\alpha\ge2$.
This gives $\rho_\alpha$, which decreases in $\alpha$; the numerical bound
$\rho_2<0.0038$ is an exact rational comparison. Neumann's series bounds
$\|(I+\Lambda^{-1}E\Lambda)^{-1}\|$ by $1/(1-\rho_2)<1.004$, and
$N^{-1}=\Lambda(I+\Lambda^{-1}E\Lambda)^{-1}\Lambda^{-1}$. For the tail,
$k(k-\frac32)\ge(m+1)(m-\frac12)\ge m^2$ for $k\ge m+1\ge2$, so the terms with
$|h|>m$ total less than $2.01e^{-\pi\alpha^2m^2}$.
\end{proof}

We now use an integer $\alpha$ with $2\le\alpha<\sqrt p$, where $p>100$ and
$t<2^p$. Put
\[
 L=\Bigl\lceil\frac{4533\,\alpha^2s}{16000\,(t-s)}\Bigr\rceil,\qquad
 w=3p+2L+12,\qquad m=\Bigl\lceil\frac{\sqrt w}{\alpha}\Bigr\rceil,
\]
\[
 \mathcal B=\{s-m,\ldots,s-1\},\qquad
 J_k=\bigl((k,k+m]\cup\mathcal B\bigr)\cap(k,s),
\]
\[
 \mathcal P=\{(i,j):|i-j|\le m\}\cup(\mathcal B\times[0,s))\cup([0,s)\times\mathcal B),
\]
with indices in $[0,s)$, and assume $L\le p$ and $s\ge4m$. Since
$\pi\log_2e<4.533$ and $s/(t-s)=1/\theta$, we have $\log_2R\le L$. Write
$Q_w$ for truncation toward zero of the real and imaginary parts to
$2^{-w}\mathbb Z$, so that $|Q_wz-z|<\sqrt2\,2^{-w}$ and $|Q_wz|\le|z|$.

\begin{lemma}[Precomputed banded solve]\label{lem:banded-solve}
Under these hypotheses put $\mathsf J=N^{-1}/2^{L+2}$, so that
$\|\mathsf J\|\le0.251$.
\begin{enumerate}
\item A fixed multitape algorithm, given $p,s,t,\alpha$, writes a factor
stream of $O(smw)$ bits in time $O(sm^2w^2)$.
\item A fixed multitape algorithm, given that stream and $u\in\mathbb D_p^s$ in
the scalar component format and increasing coordinate order, returns
$\widetilde{\mathsf J}u\in\mathbb D_p^s$ in the same format with
$2^p\|\widetilde{\mathsf J}u-\mathsf Ju\|<2$, in time $O(smw^{1+\delta})$.
\end{enumerate}
The tape count and the implicit constants are independent of $p,s,t,\alpha$.
\end{lemma}

\begin{proof}
Write $X_\Lambda=\Lambda^{-1}X\Lambda$ for matrices and $y_\Lambda=\Lambda^{-1}y$
for vectors. The machine never uses $\Lambda$; it appears only in the
analysis. As $\Lambda$ has diagonal entries in $[1,R]$ and $R\le2^L$, an
entrywise perturbation of size $\eta$ has scaled entries at most $2^L\eta$, and
$\|y_\Lambda\|\le\|y\|\le2^L\|y_\Lambda\|$.

\emph{Truncation.} Let $E_m$ keep exactly the terms with $1\le|h|\le m$. Since
$s\ge4m$, these $h$ give distinct columns different from $\ell$, so $E_m$ is
supported in the cyclic band; a band pair that wraps past $s$ has its row in
$\mathcal B$, and one that wraps past $0$ has its column in $\mathcal B$, so
the band lies in $\mathcal P$. By Lemma~\ref{lem:scaled-dominance},
$(E_m)_\Lambda$ has row sums at most $\rho_2$, and as $\alpha^2m^2\ge w$ and
$\pi\log_2e>1$, $\|(E-E_m)_\Lambda\|<3e^{-\pi w}<3\cdot2^{-w}$.

\emph{Entries.} For $1\le|h|\le m$ the exponent of a term of $E_m$ is
$\pi j/k$ with $k=s^2$ and the positive integer
$j=\alpha^2\bigl((th+\sigma_\ell)^2-\sigma_{\ell+h}^2\bigr)$, where
$\sigma_\ell=s\beta_\ell=t\ell-sq_\ell$ is an integer of $O(p)$ bits, and
$j>0$ by part (c) of the proof of Lemma~\ref{lem:scaled-dominance}.
\cite[Lemma~2.13]{HarveyHoeven2021}, applied with precision
parameter $w$ (its Section~2.2 fixes an arbitrary precision parameter larger
than $100$), returns an approximation of the real number $x$ on the grid
$2^{-w}\mathbb Z[i]$ with error less than $2\cdot2^{-w}$ in time
$O(w^{1+\delta})$; its real part $\tilde x\in2^{-w}\mathbb Z$ satisfies
$|\tilde x-x|<2\cdot2^{-w}$. Let
$\widetilde A=I+\widetilde E_m$ be the matrix of these values. Then
$\|(\widetilde E_m-E_m)_\Lambda\|<4m2^{L-w}$.

\emph{Elimination.} Put $a^{(0)}=\widetilde A$. For $k=0,\ldots,s-1$ let
$\pi_k=a^{(k)}_{kk}$, $\tilde u_{kj}=a^{(k)}_{kj}$ for $j\in J_k$,
$\psi_k=Q_w(1/\pi_k)$ and $\tilde l_{ik}=Q_w(a^{(k)}_{ik}/\pi_k)$ for
$i\in J_k$, and for $i,j\in J_k$ set
$a^{(k+1)}_{ij}=Q_w(a^{(k)}_{ij}-\tilde l_{ik}\tilde u_{kj})$, the product and
difference being exact; leave all other entries unchanged. Put
$\tilde l_{kk}=1$, $\tilde u_{kk}=\pi_k$, and let every other factor entry be
zero. If $(i,k),(k,j)\in\mathcal P$ and $i,j>k$, then $(i,j)\in\mathcal P$:
if $k\in\mathcal B$ then $i,j\in\mathcal B$, and otherwise $i$ or $j$ lies in
$\mathcal B$, or $i,j\in(k,k+m]$. Thus $J_k$ lists exactly the possible
nonzero positions, and nothing outside $\mathcal P$ is ever written.

Let $\varepsilon^{(r)}_{ij}$ be the truncation error of the update of $(i,j)$
at step $r$, or zero, and for $i>j$ let
$\varepsilon'_{ij}=\pi_j\tilde l_{ij}-a^{(j)}_{ij}$. Then
$|\varepsilon^{(r)}_{ij}|<2^{-w}$ and $|\varepsilon'_{ij}|<|\pi_j|2^{-w}$, and
telescoping the updates gives
\[
 \widetilde L\widetilde U=\widetilde A+\Delta,\qquad
 \Delta_{ij}=\sum_{r<\min(i,j)}\varepsilon^{(r)}_{ij}+[i>j]\,\varepsilon'_{ij},
\]
with $\widetilde L=(\tilde l_{ij})$ and $\widetilde U=(\tilde u_{ij})$. For
$k\le s$ let $\Delta^{(k)}$ keep only the terms with $r<k$ and the
$\varepsilon'_{ij}$ with $j<k$, and put $M^{(k)}=\widetilde A+\Delta^{(k)}$. The
same telescoping shows that the computed $\tilde l_{ir},\tilde u_{rj}$,
$r<k$, satisfy the Crout equations of $M^{(k)}$ in all positions with
$\min(i,j)<k$, and that $(a^{(k)}_{ij})_{i,j\ge k}$ is the exact $k$-th Schur
complement of $M^{(k)}$.

Suppose $|\pi_r|\le\frac{11}{10}$ for $r<k$. Then $|\Delta^{(k)}_{ij}|<(s+1)2^{-w}$,
and since a row has at most $s$ entries and $s<2^p$,
$\|\Delta^{(k)}_\Lambda\|<2^Ls(s+1)2^{-w}<2^{2p+L-w}=2^{-p-L-12}$. Hence
$M^{(k)}_\Lambda$ satisfies Lemma~\ref{lem:dominant-lu} with
$\rho\le\rho_2+4m2^{L-w}+2^{-p-L-12}<\frac1{10}$. Pivots are invariant under the
diagonal similarity, so $\pi_k\in[\frac9{10},\frac{11}{10}]$; the unscaled
active entries satisfy $|a^{(k)}_{ij}|\le\frac{11}{10}2^L$ and
$|\tilde l_{ik}|\le2^L/9+2^{-w}$. This closes the induction. All stored values
therefore fit in words of $L+3$ integer and sign bits and $w$ fractional bits.
At the end, $(\widetilde L_\Lambda,\widetilde U_\Lambda)$ is the LU factorization
of $M_\Lambda$, where $M=\widetilde A+\Delta$, and it satisfies the bounds of
Lemma~\ref{lem:dominant-lu}. Also $|\psi_k-1/\pi_k|<2^{-w}$.

\emph{Factor stream and precomputation cost.} The record $F_k$ consists of
$\psi_k$, the values $\tilde u_{kj}$ for $j\in J_k$, and the values
$\tilde l_{ik}$ for $i\in J_k$, in increasing index order and followed by a
separator: $O(m)$ words, hence $O(smw)$ bits in all. At step $k$ only rows and
columns in $W_k=\{k\}\cup J_k$ are read or written. An active entry outside
$W_k\times W_k$ is an unchanged entry of $\widetilde A$, because a row or
column index larger than $k+m$ and outside $\mathcal B$ is in
$\mathcal P$-relation with no earlier pivot. Keep the at most $(2m+1)^2$ words
of $W_k\times W_k$ row-major on a work tape. A step performs at most
$(2m+1)^2$ products and divisions of $O(w)$-bit numbers by schoolbook
arithmetic, rewrites the window once, appends the new row and column with the
$O(m)$ new entries of $\widetilde A$, and writes $F_k$. The integers
$\sigma_\ell$ come from the counters $tj=sf_j+e_j$ of the selected coordinates,
since $\sigma_\ell=e_\ell-s\,\mathbf 1_{2e_\ell\ge s}$. This proves the first
claim.

\emph{Solve.} Convert the input records to the fixed-point words and put
$v=u/2$, which is exact. Forward: for $k=0,\ldots,s-1$ let
$\tilde y_k=Q_w(\mathrm{acc}_k)$, where
$\mathrm{acc}_k=v_k-\sum_{r<k}\tilde l_{kr}\tilde y_r$ is accumulated exactly.
Keep the exact accumulators of the rows in $(k,k+m]\setminus\mathcal B$ in a
queue on one tape and those of the rows in $\mathcal B$ on another, initialized
with the corresponding $v_i$; after $\tilde y_k$ is formed, one forward scan of
$F_k$ alongside both tapes subtracts $\tilde l_{ik}\tilde y_k$ for every
$i\in J_k$, the head row leaves the queue, and the next row enters it. Backward:
for $i=s-1,\ldots,0$ form exactly
$g_i=\tilde y_i-\sum_{j\in J_i}\tilde u_{ij}\tilde x_j$ and
$\tilde x_i=Q_w(\psi_ig_i)$, reading $F_i$ backward record by record; the
values $\tilde x_j$ with $j\in(i,i+m]\setminus\mathcal B$ are kept in a queue,
and those with $j\in\mathcal B$, computed first, on a separate tape. Output
$\widetilde{\mathsf J}u=Q_p(\tilde x/2^{L+1})$ in the scalar component format.

\emph{Errors.} Forward, $\widetilde L\tilde y=v+\xi$ with
$|\xi_k|<\sqrt2\,2^{-w}$. Scaling gives
$\widetilde L_\Lambda\tilde y_\Lambda=v_\Lambda+\xi_\Lambda$ with
$\|\xi_\Lambda\|<\sqrt2\,2^{-w}$ and $\|v_\Lambda\|\le\frac12$, hence
$\|\tilde y_\Lambda-\widetilde L_\Lambda^{-1}v_\Lambda\|\le\frac87\sqrt2\,2^{-w}$ and
$\|\tilde y_\Lambda\|<0.58$. Backward, downward induction gives
$|\tilde x_j|\le\lambda_j$: if this holds for $j>i$, then
$|g_i|\le(0.58+\frac1{10})\lambda_i$ and $|\tilde x_i|\le
(\frac{10}9+2^{-w})0.68\lambda_i<\lambda_i$. The residual
$\eta_i=\tilde x_i-g_i/\pi_i$ satisfies
$|\eta_i|/\lambda_i<(\sqrt2+0.68)2^{-w}<2.2\cdot2^{-w}$, and
$\widetilde U\tilde x=\tilde y+D_U\eta$. Lemma~\ref{lem:dominant-lu} for the
scaled factors gives
\[
 \|\tilde x_\Lambda-M_\Lambda^{-1}v_\Lambda\|
 \le\tfrac54\cdot\tfrac87\sqrt2\,2^{-w}+\tfrac98\cdot2.2\cdot2^{-w}<5\cdot2^{-w},
\]
hence $\|\tilde x-M^{-1}v\|<5\cdot2^{L-w}$. Further,
$\|N_\Lambda-M_\Lambda\|<3\cdot2^{-w}+4m2^{L-w}+2^{-p-L-12}<2^{-p-L-11}$,
$\|N_\Lambda^{-1}\|\le1.004$ and $\|M_\Lambda^{-1}\|\le\frac{10}9$, so
\[
 \|(N^{-1}-M^{-1})v\|\le2^L\cdot1.004\cdot\tfrac{10}9\cdot2^{-p-L-11}\cdot\tfrac12
 <0.56\cdot2^{-p-11},
\]
and $\|\tilde x-N^{-1}v\|<2^{-p-10}$. Since $N^{-1}v/2^{L+1}=\mathsf Ju$, the
output error is less than $\sqrt2+2^{-11}<2$ units of $2^{-p}$. Also
$\|\mathsf J\|\le1.004R/2^{L+2}\le0.251$, so the truncated output has modulus
less than one and lies in $\mathbb D_p^s$.

Every exact accumulator is bounded by $\frac12+0.58\cdot2^L/8$ and every $g_i$
by $0.68\cdot2^L$, so all words have $O(w)$ bits; as $L\le p$, $w=O(p)$.

\emph{Solve cost.} Each index uses $O(m)$ products of $O(w)$-bit integers,
$O(w^{1+\delta})$ each with the established multiplier, and $O(mw)$ head moves
for one record of the factor stream and one pass over each queue or the
$\mathcal B$ tape. The factor stream is traversed once forward and once
backward. Initializing the $\mathcal B$ accumulators, reversing the output
order and the record conversions take $O(sw)$. Hence the solve costs
$O(smw^{1+\delta})$ on a fixed number of tapes.
\end{proof}


\subsection{The revised one-dimensional interface}
\begin{proposition}\label{prop:note-one-dimensional}
Fix $0<\delta<1/8$. Let $p>100$, $2\le s<t<2^p$ and $2\le\alpha<\sqrt p$ be
integers with $\gcd(s,t)=1$ and $t<2s$, let $L,w,m$ be as in
Lemma~\ref{lem:banded-solve}, and assume $L\le p$ and $s\ge4m$. With
$(P_su)_j=u_{tj}$ and $(P_tv)_k=v_{-sk}$, the maps $A=\mathsf S/2$ and
$B_0=\mathsf D'\mathsf JC$, where $\mathsf D'=2^{-(2\alpha^2-2)}\mathsf D$,
are contractions and
\[
 P_sF_s=2^{2\alpha^2+L+1}B_0P_tF_tA .
\]
There are fixed multitape algorithms for approximations $\widetilde A$ and
$\widetilde B_0$ with errors less than $p^2$ and times
$O(tp^{3/2+\delta}\alpha)$ and $O(tmw^{1+\delta}+tp^{1+\delta})$, the second
given the factor stream of Lemma~\ref{lem:banded-solve}.
\end{proposition}
\begin{proof}
By Lemma~\ref{lem:scaled-dominance} the matrix $N$ is invertible, so the
left-inverse identity $(\mathsf DN^{-1}C)\mathsf T=I$
of~\cite[(4.2)]{HarveyHoeven2021} holds, and the resampling
identity~\cite[Theorem~4.2]{HarveyHoeven2021} gives
$P_sF_s=\mathsf DN^{-1}CP_tF_t\mathsf S$. Inserting
$\mathsf D=2^{2\alpha^2-2}\mathsf D'$, $N^{-1}=2^{L+2}\mathsf J$ and
$\mathsf S=2A$ proves the identity. The norms $\|A\|<\frac34$ and
$\|\mathsf D'\|<1$ hold as in the proof of~\cite[Proposition~4.7(i)]{HarveyHoeven2021},
and $\|\mathsf J\|\le0.251$ by Lemma~\ref{lem:banded-solve}. The
approximations are $\widetilde A=\widetilde{\mathsf S}'$ and
$\widetilde B_0=\widetilde{\mathsf D}'\widetilde{\mathsf J}C$, with the cited
algorithms for $\mathsf S'$ (error $<16p$) and $\mathsf D'$ (error $<4$), the
banded solve (error $<2$), and the exact ordered selection $C$. Errors add
through the contractions, so $\varepsilon(\widetilde B_0)<4+2<p^2$ and
$\varepsilon(\widetilde A)<16p<p^2$.

The cited statements are used outside the hypothesis $\theta>p/\alpha^4$ of
their Proposition~4.7. We record the sentences of their proofs that are
used, by section and printed page of the cited manuscript; none of them
involves $\theta$.
\begin{itemize}
\item Section~4.1, p.~24: ``let $s$ and $t>s$ be positive integers such that
$\gcd(s,t)=1$, and let $\alpha\in(0,\infty)$''; Theorem~4.2, p.~24: ``We have
$\mathsf TP_sF_s=P_tF_t\mathsf S$''.
\item Lemma~4.5, p.~26: ``We have $\|\mathsf S\|<1+\alpha^{-1}$.'' The
estimate (4.1), p.~27: ``$\|\mathsf D\|\le e^{\pi\alpha^2/4}$''. The left
inverse (4.2), p.~28: ``$(\mathsf DN^{-1}C)\mathsf T=\mathsf DN^{-1}T'
=\mathsf D(T'\mathsf D)^{-1}T'=I$'', valid whenever $N=T'\mathsf D$ is
invertible.
\item Proof of Proposition~4.7(i), p.~29: ``Lemma~4.5 implies that
$\|\mathsf S\|<\frac32$'' and ``By (4.1) we have
$\|\mathsf D\|\le e^{\pi\alpha^2/4}<2^{1.14\alpha^2}<2^{2\alpha^2-2}$, as
$1.14x<2x-2$ for all $x>4$; hence also $\|\mathsf D'\|<1$''. The last
inequality holds for every $x>2.33$, in particular for $x=\alpha^2=4$.
\item Lemma~4.8, pp.~29--30, computing $\mathsf D'$: its proof uses that
``$\alpha$, $s$, $t$ and $\ell$ are all integers with $O(p)$ bits (here we
have used the hypotheses $s,t<2^p$ and $\alpha<p^{1/2}$)'', the estimate
$\|\mathsf D'\|<1$, and Lemma~2.14 ``with $\sigma:=2\alpha^2-2<2p$''.
\item Lemma~4.9, pp.~30--31, computing $\mathsf S'$ with $m=\lceil p^{1/2}\rceil\alpha$
terms: its proof uses ``Since $\|u\|\le1$ and $\alpha>2$ we have'', only to
bound the weight $\frac12\alpha^{-1}$ by $\frac14$, which also holds for
$\alpha=2$; then ``Since $\alpha<p^{1/2}$ we have
$w^m=e^{-\pi\lceil p^{1/2}\rceil/\alpha}\le e^{-\pi p^{1/2}/\alpha}<e^{-\pi}<0.05$'',
``As $m<\lceil p^{1/2}\rceil p^{1/2}\le p+p^{1/2}$ and $p>100$'', and
``Recalling from the proof of Proposition~4.7(i) that $\|\mathsf S'\|<\frac34$''.
Its statement assumes the hypotheses of Proposition~4.7, which allow
$2\le\alpha<p^{1/2}$.
\item Section~2.2, p.~8: ``Fix a precision parameter $p>100$''; Lemma~2.13,
p.~17, is stated for this arbitrary precision and is applied above with
precision $w$.
\end{itemize}
The cited Lemmas~4.11 and~4.12 (the Neumann solve), and Lemma~4.6, are not
used.
\end{proof}

For the tensor interface, every axis $i$ uses its own $L_i$, the scale
becomes $\gamma=\sum_i(2\alpha^2+L_i+1)$, and the factor stream of axis $i$
is written once on a dedicated tape before its lines are compressed. The
error bound $dp^2$ and the contraction properties are unchanged. The
one-dimensional calls cost $\sum_i(T/t_i)O(t_ip^{3/2+\delta}\alpha
+t_im_iw_i^{1+\delta})$, and the precomputations
$O(\sum_is_im_i^2w_i^2)$.

\section{Parameters, costs and the witness}
Keep $b=\lceil\log_2n\rceil$, $p=6b$, $d=\lfloor b^\epsilon\rfloor$, $K=d$,
$T,r,\ell,t_i$, $\eta=1/(4d)$ and the prime selection
$(1-2\eta)t_i<s_i\le(1-\eta)t_i$ of the pinned source. Set
\[
 \alpha=2,\qquad
 L_i=\Bigl\lceil\frac{4533\,s_i}{4000\,(t_i-s_i)}\Bigr\rceil,\qquad
 \gamma=\sum_{i=1}^d(2\alpha^2+L_i+1).
\]
Then $\theta_i>\eta$ gives $s_i/(t_i-s_i)<4d$, so $L_i\le\lceil4.533d\rceil$,
and for $d\ge22$
\[
 \gamma\le d(10+4.533d)\le5d^2\le5b^{2\epsilon}.
\]
Also $s_i>(1-2\eta)t_i\ge t_i/2$, $L_i\le p$, $w_i\le6p$,
$m_i\le\sqrt{6p}$ and $s_i\ge r/4\ge4m_i$ for large $n$, so
Proposition~\ref{prop:note-one-dimensional} applies on every axis. Every
later use of $\gamma$ in the pinned precision analysis needs only
$\gamma\le b/4$, which holds once $b^{1-2\epsilon}\ge20$. The Gaussian line
maps cost $O(dTp^{3/2+\delta})$ in total, and the precomputations
$O(drp^3)=o(Tp)$ because $T/r\ge(r/2)^{d-1}$ and $r$ exceeds every fixed
power of $p$. After division by $Tp$, the Gaussian row has power
$\frac12+\delta+\epsilon$. No other row of the cost table changes. The seven
margins are
\begin{align*}
 g_1&=1-\epsilon(1+c), &g_2&=\epsilon c\,\ab,\\
 g_3&=\epsilon(1-\lambda'), &g_4&=\ab(1-\epsilon),\\
 g_5&=\tfrac12-\delta-\epsilon,
 &g_6&=1-\delta-\epsilon,\qquad g_7=\epsilon.
\end{align*}
Relative to the baseline constraint list, exactly five entries change:
the guard $1-\epsilon(5-4\beta+\zeta)$ becomes $1-\epsilon$; the Gaussian
cost $\frac14-\delta-\frac54\epsilon$ becomes $\frac12-\delta-\epsilon$; the
dimension and source-scale conditions $\frac13-\epsilon$ and
$\frac12-\frac32\epsilon$ become $\frac12-\epsilon$ and $1-2\epsilon$; and
$\alpha<\sqrt p$ becomes unconditional.

Put $t=10^{-12}$ and choose
\[
\begin{gathered}
 \tau=1-\ab,\quad\sigma=1-\ac,\quad\beta=\frac1{10},\quad c=1,\quad
 \delta=\frac1{10^{12}},\quad
 \epsilon=\frac{124999999}{250000000},\\
 \lambda=1-\ab(1-t/2),\qquad\lambda'=1-\ab(1-t),\qquad
 \kappa=\frac{296461013}{2\cdot10^{17}}.
\end{gathered}
\]
Since $\sigma<\tau$, the compact-control recurrence has
$\chi=\tau+(1-\beta)\max\{\sigma-\tau,0\}=\tau$, and exact comparisons give
$\max\{\tau,\sigma,\chi\}<\lambda<\lambda'<1$,
$\sigma+\beta(1-\sigma)<\lambda'$ and $1-c=0<\lambda'$. All the other strict
conditions hold, with $1-2\epsilon=8/10^9$. The minimum margin is
\[
\begin{gathered}
 G_*=g_3=\epsilon\ab(1-t)=\frac{3705762684806317521666189976715649}{25\cdot10^{41}},\\
 G_*-\kappa=\frac{22306317521666189976715649}{25\cdot10^{41}}>0,
\end{gathered}
\]
while $g_5=3999/10^{12}$, $g_1=8/10^9$ and $g_2,g_4>G_*$. The strict gap
absorbs all fixed powers of $\log p$.

\begin{theorem}[Conditional refinement]
Under the retained upstream, network and compact-control interfaces, and the
written proofs of this note,
\[
 T(n)=O\!\left(n(\log n)^{1-296461013/(2\cdot10^{17})}\right).
\]
The saving lies strictly between $2^{-30}$ and $2^{-29}$ and exceeds the
previous witness $5929220328/10^{19}$ by the factor
$7411525325/2964610164\approx2.49999997$.
\end{theorem}

The choice of $\epsilon$ is within $4\cdot10^{-9}$ of $\frac12$. The resulting
eventual cutoffs, for example $b\ge2^{541000000}$ for $\gamma\le b/4$, are
astronomically large; they are fixed constants, which the asymptotic
statement allows. Any fixed $\epsilon<\frac12$ gives a witness with smaller
cutoffs; $\epsilon=\frac25$ gives about twice the previous witness with
$b^{1/5}\ge20$.

\section{The supremum and the ceiling of this assembly}
Write $q=1-\lambda'$. For the current savings ($\sigma<\tau$) the layer
conditions can be met exactly when $q<\min(\ab,(1-\beta)\ac)$, and $\beta$
now enters no other condition. With
$c=1$, the binding margins are $g_3=\epsilon q$ and
$g_5=\frac12-\delta-\epsilon$; their crossing $\epsilon\to1/(2+2q)$ lies
strictly below every other cap. Hence the parameter supremum is
\[
 \kappa_{\sup}=\frac Q{2+2Q},\qquad Q=\min(\ab,\ac),
\]
about $1.4823050814\cdot10^{-9}$ for the current savings. The previous
assembly gives $Q/(5+4Q)$.

The value is within the factor $1+Q$ of a ceiling that no Gaussian stage
can remove. The retained conditions include $g_2=\epsilon c\ab>\kappa$,
$g_3=\epsilon q>\kappa$, the transform layout requirement
$K=\lfloor d^c\rfloor\le\ell-1$, which gives $\epsilon(1+c)<1$, and the
prime-interval requirement $\log(r/2)/d\to\infty$, which gives
$1-2\epsilon>0$. Therefore
\[
 \kappa<\epsilon\min(q,c\,\ab)\le\frac{\min(q,c\,\ab)}{\max(1+c,2)}
 \le\frac q2\le\frac Q2 .
\]
Exceeding $Q/2$ requires relaxing both $\epsilon<\frac12$ caps: the chunk
layout of the synthetic transform and the prime-interval argument. Further
gains otherwise come only from the finite networks: after the linear guard,
$\beta$ may be arbitrarily small, so improving $\ab$ alone raises $\kappa$
proportionally until $\ab$ reaches $\ac$.

\section{Reproduction and verification boundary}
The implementation uses only Python's standard library:
\begin{verbatim}
python3 scripts/assembly_lu.py
python3 scripts/make_assembly_lu_patch.py
python3 -m unittest discover -s tests -p 'test_assembly_lu*.py' -v
git apply --check --directory=upstream patches/assembly-lu-30.patch
\end{verbatim}
The certificate re-derives $\ab$ with the retained paired-bit certificate and
$\ac$ from the pair-star counts and a rational logarithm enclosure. It
recomputes the node allowance $D_{\rm node}$ and $C_0^*$ from the pair-star
circuit, the entries of $C^{\otimes f}$ for $f\le12$, rational enclosures of
$\pi$, $e^{2\pi}$, $e^{4\pi}$, $e^{14\pi}$ and $\log2$, every constant of
Lemmas~\ref{lem:dominant-lu}--\ref{lem:banded-solve}, the identities for
$\psi$ on exact samples, and both scaled-exponent inequalities exactly on
every row of six instances. It checks every strict constraint and margin of
the witness, and that the retained constraint list agrees with the
repository's code. The tests run the factorization and solve of
Lemma~\ref{lem:banded-solve} in exact arithmetic at $p=101$, including two
$\alpha=2$ instances whose unscaled off-diagonal row sums exceed $0.46$, and
compare the output with a reference solve of the full matrix at more than
$200$ digits; the worst observed error is below $1.4$ units of $2^{-p}$.

The new patch applies to the unmodified pinned manuscript. It incorporates
the earlier compact-control and pair-star modifications, replaces the guard
subsection, inserts the banded solve before the one-dimensional resampling
lemma and revises that lemma, and updates the assembly parameters, cost
table, margins and witness. The verifier does not implement the
multiplication machine, and the finite tests do not prove uniform tape
estimates or the cited analytic identities. Those remain stated proof
dependencies. No unrestricted optimality, independent review, or practical
performance claim is made.

\begin{thebibliography}{9}
\raggedright
\bibitem{fork} odepoint.
\emph{intmultspam: integer multiplication, just for fun.}
Fork of CrocSwap/integer-mult-bounds, revision \texttt{8e1701f}, October 2026.
\bibitem{baseline} Douglas Colkitt.
\emph{A sharper exponent for integer multiplication.}
CrocSwap/integer-mult-bounds, revision
\texttt{6e564879f51ae16f23d392e9e196c605f36d90df}, October 7, 2026.
\url{https://github.com/CrocSwap/integer-mult-bounds}.
\bibitem{upstream} OpenAI.
\emph{Integer multiplication below $n\log n$.}
Source pinned at revision
\texttt{adc7f1241b42e322a6451854ab7e4b4c146bf78a}; see the repository's
\texttt{upstream/manifest.json} for source URLs and hashes.
\url{https://github.com/openai/math}.
\bibitem{HarveyHoeven2021} David Harvey and Joris van der Hoeven.
\emph{Integer multiplication in time $O(n\log n)$.}
Annals of Mathematics 193 (2021). Numbering and printed pages refer to the
45-page author-hosted manuscript
\url{https://www.texmacs.org/joris/nlogn/nlogn.pdf}, as in the pinned
source; its text agrees with hal-02070778v2.
\end{thebibliography}
\end{document}
